\documentclass[11pt,a4paper]{article}
\usepackage{fontspec,amsmath,amssymb,amsthm,microtype}
\setmainfont{lmroman10-regular.otf}[BoldFont=lmroman10-bold.otf,ItalicFont=lmroman10-italic.otf,BoldItalicFont=lmroman10-bolditalic.otf]
\usepackage{xcolor}
\usepackage[most]{tcolorbox}
\usepackage{mathrsfs}
\providecommand{\lt}{<}
\providecommand{\gt}{>}
\setmonofont{DejaVuSansMono.ttf}[BoldFont=DejaVuSansMono-Bold.ttf,ItalicFont=DejaVuSansMono-Oblique.ttf,BoldItalicFont=DejaVuSansMono-BoldOblique.ttf,Scale=0.86]
\newfontfamily\leanfallback{STIXTwoMath-Regular.otf}[Scale=0.86]
\definecolor{leanrule}{gray}{0.62}
\definecolor{leanback}{gray}{0.962}
\definecolor{leantitle}{gray}{0.22}
\definecolor{leanmuted}{gray}{0.38}
\newtcolorbox{leanbox}[1]{enhanced,breakable,colback=leanback,colframe=leanrule,
  boxrule=0.4pt,arc=0.6mm,left=2.4mm,right=2.4mm,top=1.8mm,bottom=1.8mm,
  before skip=0.8em,after skip=1em,fonttitle=\small\bfseries,coltitle=leantitle,
  colbacktitle=leanback,titlerule=0.3pt,title={#1}}
\newlength{\leancharwidth}
\newenvironment{leancode}{\par\smallskip\begingroup\ttfamily\footnotesize\raggedright
  \settowidth{\leancharwidth}{M}\setlength{\parskip}{0pt}\setlength{\parindent}{0pt}
  \frenchspacing}{\par\endgroup\smallskip}
\newcommand{\leanline}[2]{\hangindent=\dimexpr#1\leancharwidth+4\leancharwidth\relax
  \hangafter=1\noindent\hspace*{#1\leancharwidth}#2\par}
\newcommand{\leansource}[1]{\par{\footnotesize\color{leanmuted}#1}\par}
\newcommand{\leanlabel}[1]{\par\smallskip\noindent{\small\bfseries #1}\par}
\usepackage[margin=25mm]{geometry}
\usepackage[unicode,hidelinks]{hyperref}
\hypersetup{pdfauthor={LangevinLean Project}}
\setlength{\emergencystretch}{3em}
\setlength{\parskip}{0.4em}
\setlength{\parindent}{0pt}
\allowdisplaybreaks[2]
\raggedbottom
\title{\texorpdfstring{Reference and auxiliary trajectories on each time interval}{Reference and auxiliary trajectories on each time interval}}
\author{LangevinLean Project}
\date{}
\begin{document}\maketitle
Draft mathematical exposition; these notes may contain mistakes.


The local KLMC analysis compares a numerical state with an exact reference state. On each step it inserts an auxiliary exact path started at the numerical state and driven by precisely the same additive input as the reference. This section constructs those paths from the proved kinetic flow. The differential, kinematic and mild identities used in the local-error theorem follow from their integral equations.



The exact reference, numerical chain, and auxiliary exact path have different initial phases but compatible local inputs. Their proved integral identities provide the complete structural foundation of the coupling argument.


\section{\texorpdfstring{Restart an integral solution at any time}{Restart an integral solution at any time}}


If \(Y_t-Y_0=-\int_0^tg(Y_s)ds+W_t\), then \(Y_{a+t}-Y_a=-\int_0^tg(Y_{a+s})ds+(W_{a+t}-W_a)\) whenever the shifted interval lies in the original horizon. The proof subtracts the two equations at \(a+t\) and \(a\), splits the continuous drift integral at \(a\), and changes variables in the remaining integral. Additive group algebra cancels the shared starting point. This supplies the exact time-shift identity required when segmenting a continuous diffusion.



Subtract the global integral equation at a restart time from its value later in the interval. The new initial phase is the global phase at the restart, and the new forcing is the corresponding input increment. Thus the shifted path solves the same kinetic equation locally.

\begin{leanbox}{Lean: \texttt{additive\_\allowbreak{}integral\_\allowbreak{}segment}}
\leansource{\texttt{additive\_\allowbreak{}integral\_\allowbreak{}segment} \quad KineticLangevin/\allowbreak{}CoupledSegments.lean, lines 12–29}
\begin{leancode}
\leanline{0}{theorem\ additive\_integral\_segment\ (g\ :\ E\ →\ E)\ (hg\ :\ Continuous\ g)}
\leanline{4}{(Y\ W\ :\ ℝ\ →\ E)\ (hY\ :\ Continuous\ Y)\ (T\ a\ h\ :\ ℝ)}
\leanline{4}{(ha\ :\ 0\ ≤\ a)\ (hh\ :\ 0\ ≤\ h)\ (hah\ :\ a+h\ ≤\ T)}
\leanline{4}{(heq\ :\ ∀\ t\ ∈\ Icc\ 0\ T,\ Y\ t\ -\ Y\ 0\ =\ -(∫\ s\ in\ 0..t,\ g\ (Y\ s))\ +\ W\ t)}
\leanline{4}{(t\ :\ ℝ)\ (ht\ :\ t\ ∈\ Icc\ 0\ h)\ :}
\leanline{4}{Y\ (a+t)\ -\ Y\ a\ =\ -(∫\ s\ in\ 0..t,\ g\ (Y\ (a+s)))\ +\ (W\ (a+t)-W\ a)}
\end{leancode}
\leanlabel{How the Lean proof follows the argument}
\texttt{additive\_\allowbreak{}integral\_\allowbreak{}segment} gives the formal counterpart of this step. The main intermediate facts are \texttt{hat}, \texttt{ha\allowbreak{}T}, \texttt{hc}, \texttt{hi}, \texttt{hs}, \texttt{hya} in that order. Named proof applications include \texttt{interval\allowbreak{}Integral.\allowbreak{}integral\_\allowbreak{}comp\_\allowbreak{}add\_\allowbreak{}left}; the adjacent proof shows their hypotheses and order of use.
\end{leanbox}

\section{\texorpdfstring{Use a raw continuous input with the selected kinetic flow}{Use a raw continuous input with the selected kinetic flow}}


A continuous input \(W\) is restricted to the compact solution horizon, giving a continuous function on that compact interval. If \(W_0=0\), the flow's normalized clipped extension equals \(W_t\) at every time in that horizon. The interval projection fixes every point in the solution horizon, so clipping does not change the input there. The selected segment is then the kinetic solution map given by Picard existence with that compact input. Its specification gives continuity, its initial condition, and the exact integral equation with the original raw input.



The selected flow only needs a continuous additive velocity input starting at zero. Subtracting the input’s initial value supplies that normalization for a raw continuous path. This permits local segments to be defined uniformly without choosing a new stochastic solution for each interval.

\begin{leanbox}{Lean: \texttt{compactVelocityInput}, \texttt{normalizedForcing\_\allowbreak{}compactVelocityInput}, \texttt{kineticSegmentFlow} and 1 more}
\leansource{\texttt{compactVelocityInput} \quad KineticLangevin/\allowbreak{}CoupledSegments.lean, lines 32–33}
\begin{leancode}
\leanline{0}{def\ compactVelocityInput\ (T\ :\ ℝ)\ (W\ :\ ℝ\ →\ E)\ (hW\ :\ Continuous\ W)\ :\ VelocityForcing\ (E\ :=\ E)\ T\ :=}
\leanline{2}{⟨fun\ t\ =>\ W\ t,hW.comp\ continuous\_subtype\_val⟩}
\end{leancode}
\leansource{\texttt{normalizedForcing\_\allowbreak{}compactVelocityInput} \quad KineticLangevin/\allowbreak{}CoupledSegments.lean, lines 36–40}
\begin{leancode}
\leanline{0}{theorem\ normalizedForcing\_compactVelocityInput\ (T\ :\ ℝ)\ (hT\ :\ 0\ ≤\ T)}
\leanline{4}{(W\ :\ ℝ\ →\ E)\ (hW\ :\ Continuous\ W)\ (hW0\ :\ W\ 0\ =\ 0)\ (t\ :\ ℝ)\ (ht\ :\ t\ ∈\ Icc\ 0\ T)\ :}
\leanline{4}{normalizedForcing\ T\ hT\ (compactVelocityInput\ T\ W\ hW)\ t\ =\ W\ t}
\end{leancode}
\leansource{\texttt{kineticSegmentFlow} \quad KineticLangevin/\allowbreak{}CoupledSegments.lean, lines 43–45}
\begin{leancode}
\leanline{0}{def\ kineticSegmentFlow\ (g\ :\ E\ →\ E)\ (L\ γ\ u\ :\ ℝ≥0)\ (hg\ :\ LipschitzWith\ L\ g)}
\leanline{4}{(T\ :\ ℝ)\ (hT\ :\ 0\ ≤\ T)\ (initial\ :\ E\ ×\ E)\ (W\ :\ ℝ\ →\ E)\ (hW\ :\ Continuous\ W)\ :\ ℝ\ →\ E\ ×\ E\ :=}
\leanline{2}{kineticFlow\ g\ L\ γ\ u\ hg\ T\ hT\ initial\ (compactVelocityInput\ T\ W\ hW)}
\end{leancode}
\leansource{\texttt{kineticSegmentFlow\_\allowbreak{}spec} \quad KineticLangevin/\allowbreak{}CoupledSegments.lean, lines 48–60}
\begin{leancode}
\leanline{0}{theorem\ kineticSegmentFlow\_spec\ (g\ :\ E\ →\ E)\ (L\ γ\ u\ :\ ℝ≥0)\ (hg\ :\ LipschitzWith\ L\ g)}
\leanline{4}{(T\ :\ ℝ)\ (hT\ :\ 0\ ≤\ T)\ (initial\ :\ E\ ×\ E)\ (W\ :\ ℝ\ →\ E)\ (hW\ :\ Continuous\ W)}
\leanline{4}{(hW0\ :\ W\ 0\ =\ 0)\ :}
\leanline{4}{Continuous\ (kineticSegmentFlow\ g\ L\ γ\ u\ hg\ T\ hT\ initial\ W\ hW)\ ∧}
\leanline{6}{kineticSegmentFlow\ g\ L\ γ\ u\ hg\ T\ hT\ initial\ W\ hW\ 0\ =\ initial\ ∧}
\leanline{6}{∀\ t\ ∈\ Icc\ 0\ T,}
\leanline{8}{kineticSegmentFlow\ g\ L\ γ\ u\ hg\ T\ hT\ initial\ W\ hW\ t\ -\ initial\ =}
\leanline{10}{-(∫\ s\ in\ 0..t,\ kineticNegativeDrift\ g\ γ\ u}
\leanline{12}{(kineticSegmentFlow\ g\ L\ γ\ u\ hg\ T\ hT\ initial\ W\ hW\ s))\ +\ (W\ t,0)}
\end{leancode}
\leanlabel{How the Lean proof follows the argument}
\texttt{compact\allowbreak{}Velocity\allowbreak{}Input} names the mathematical object defined above. Its defining expression appears in the adjacent statement; subsequent identities may unfold this definition directly.




\texttt{normalized\allowbreak{}Forcing\_\allowbreak{}compact\allowbreak{}Velocity\allowbreak{}Input} gives the formal counterpart of this step.




\texttt{kinetic\allowbreak{}Segment\allowbreak{}Flow} names the mathematical object defined above. Its defining expression appears in the adjacent statement; subsequent identities may unfold this definition directly.




\texttt{kinetic\allowbreak{}Segment\allowbreak{}Flow\_\allowbreak{}spec} gives the formal counterpart of this step. The main intermediate facts are \texttt{hs} in that order. Named proof applications include \texttt{hs.\allowbreak{}2.\allowbreak{}2}; the adjacent proof shows their hypotheses and order of use.
\end{leanbox}

\section{\texorpdfstring{Derive velocity, kinematic and mild equations by projection}{Derive velocity, kinematic and mild equations by projection}}


Applying the velocity continuous linear projection to the kinetic integral equation gives \(V_t-V_0=-\int_0^t(\gamma V_s+u g(X_s))ds+W_t\). Applying the position projection gives \(X_t-X_0=\int_0^tV_sds\), because the positional input is exactly zero. In both cases the continuous linear projection commutes with the Bochner integral. If the filtered input satisfies \(B_t=W_t-\gamma\int_0^tB_sds\), the proved Banach-space integrating-factor theorem gives \(V_T=\psi_0(T)V_0-\int_0^Te^{-\gamma(T-s)}u g(X_s)ds+B_T\). Thus the exponential mild equation is a consequence of the kinetic equation, rather than a new stochastic hypothesis.



Projection of the phase equation gives the velocity and position integral identities. The integrating-factor argument then gives the mild velocity formula, while the position equation is already the kinematic primitive. These are the exact analytic identities required by the local-error estimate.

\begin{leanbox}{Lean: \texttt{kinetic\_\allowbreak{}integral\_\allowbreak{}velocity}, \texttt{kinetic\_\allowbreak{}integral\_\allowbreak{}kinematic}, \texttt{kinetic\_\allowbreak{}integral\_\allowbreak{}mild}}
\leansource{\texttt{kinetic\_\allowbreak{}integral\_\allowbreak{}velocity} \quad KineticLangevin/\allowbreak{}CoupledSegments.lean, lines 63–74}
\begin{leancode}
\leanline{0}{theorem\ kinetic\_integral\_velocity\ (g\ :\ E\ →\ E)\ (γ\ u\ T\ :\ ℝ)\ (Y\ :\ ℝ\ →\ E\ ×\ E)}
\leanline{4}{(W\ :\ ℝ\ →\ E)\ (hY\ :\ Continuous\ Y)\ (hg\ :\ Continuous\ g)}
\leanline{4}{(heq\ :\ ∀\ t\ ∈\ Icc\ 0\ T,\ Y\ t\ -\ Y\ 0\ =}
\leanline{6}{-(∫\ s\ in\ 0..t,\ kineticNegativeDrift\ g\ γ\ u\ (Y\ s))\ +\ (W\ t,0))}
\leanline{4}{(t\ :\ ℝ)\ (ht\ :\ t\ ∈\ Icc\ 0\ T)\ :}
\leanline{4}{(Y\ t).1\ -\ (Y\ 0).1\ =\ -(∫\ s\ in\ 0..t,\ γ\ •\ (Y\ s).1\ +\ u\ •\ g\ (Y\ s).2)\ +\ W\ t}
\end{leancode}
\leansource{\texttt{kinetic\_\allowbreak{}integral\_\allowbreak{}kinematic} \quad KineticLangevin/\allowbreak{}CoupledSegments.lean, lines 77–89}
\begin{leancode}
\leanline{0}{theorem\ kinetic\_integral\_kinematic\ (g\ :\ E\ →\ E)\ (γ\ u\ T\ :\ ℝ)\ (Y\ :\ ℝ\ →\ E\ ×\ E)}
\leanline{4}{(W\ :\ ℝ\ →\ E)\ (hY\ :\ Continuous\ Y)\ (hg\ :\ Continuous\ g)}
\leanline{4}{(heq\ :\ ∀\ t\ ∈\ Icc\ 0\ T,\ Y\ t\ -\ Y\ 0\ =}
\leanline{6}{-(∫\ s\ in\ 0..t,\ kineticNegativeDrift\ g\ γ\ u\ (Y\ s))\ +\ (W\ t,0))}
\leanline{4}{(t\ :\ ℝ)\ (ht\ :\ t\ ∈\ Icc\ 0\ T)\ :}
\leanline{4}{(Y\ t).2\ -\ (Y\ 0).2\ =\ ∫\ s\ in\ 0..t,\ (Y\ s).1}
\end{leancode}
\leansource{\texttt{kinetic\_\allowbreak{}integral\_\allowbreak{}mild} \quad KineticLangevin/\allowbreak{}CoupledSegments.lean, lines 92–102}
\begin{leancode}
\leanline{0}{theorem\ kinetic\_integral\_mild\ (g\ :\ E\ →\ E)\ (γ\ u\ T\ :\ ℝ)\ (hT\ :\ 0\ ≤\ T)}
\leanline{4}{(Y\ :\ ℝ\ →\ E\ ×\ E)\ (W\ B\ :\ ℝ\ →\ E)\ (hY\ :\ Continuous\ Y)\ (hg\ :\ Continuous\ g)}
\leanline{4}{(hB\ :\ Continuous\ B)}
\leanline{4}{(heq\ :\ ∀\ t\ ∈\ Icc\ 0\ T,\ Y\ t\ -\ Y\ 0\ =}
\leanline{6}{-(∫\ s\ in\ 0..t,\ kineticNegativeDrift\ g\ γ\ u\ (Y\ s))\ +\ (W\ t,0))}
\leanline{4}{(hnoise\ :\ ∀\ t\ ∈\ Icc\ 0\ T,\ B\ t\ =\ W\ t\ -\ γ\ •\ ∫\ s\ in\ 0..t,\ B\ s)\ :}
\leanline{4}{(Y\ T).1\ =\ psi\ γ\ 0\ T\ •\ (Y\ 0).1\ -}
\leanline{6}{(∫\ s\ in\ 0..T,\ Real.exp\ (-γ*(T-s))\ •\ (u\ •\ g\ (Y\ s).2))\ +\ B\ T}
\end{leancode}
\leanlabel{How the Lean proof follows the argument}
\texttt{kinetic\_\allowbreak{}integral\_\allowbreak{}velocity} gives the formal counterpart of this step. The main intermediate facts are \texttt{hc}, \texttt{hp} in that order. Named proof applications include \texttt{hp}; the adjacent proof shows their hypotheses and order of use.




\texttt{kinetic\_\allowbreak{}integral\_\allowbreak{}kinematic} gives the formal counterpart of this step. The main intermediate facts are \texttt{hc}, \texttt{hp} in that order. Named proof applications include \texttt{hp}; the adjacent proof shows their hypotheses and order of use.




\texttt{kinetic\_\allowbreak{}integral\_\allowbreak{}mild} gives the formal counterpart of this step. Named proof applications include \texttt{damped\_\allowbreak{}integral\_\allowbreak{}eq\_\allowbreak{}mild}; the adjacent proof shows their hypotheses and order of use.
\end{leanbox}

\section{\texorpdfstring{Construct exact reference segments and the defined numerical chain}{Construct exact reference segments and the defined numerical chain}}


The reference chain starts from its supplied phase vector. At each iteration, its next state is the endpoint of the selected exact kinetic segment driven by that iteration's input. The corresponding reference segment is the same selected path, so its endpoint matches the next reference state by definition. The numerical chain instead performs the frozen-force exponential update with innovation \((B_n(h),\int_0^hB_n(s)ds)\). This uses the constructed filtered input and its positional primitive.



The reference uses exact selected segments; the numerical sequence uses the defined exponential update with the same interval responses. Their endpoint matching is built into the recursions. The comparison is therefore between trajectories rather than arbitrary sequences satisfying a desired bound.

\begin{leanbox}{Lean: \texttt{kineticReferenceChain}, \texttt{kineticReferenceSegment}, \texttt{exponentialInputChain} and 1 more}
\leansource{\texttt{kineticReferenceChain} \quad KineticLangevin/\allowbreak{}CoupledSegments.lean, lines 105–110}
\begin{leancode}
\leanline{0}{def\ kineticReferenceChain\ (g\ :\ E\ →\ E)\ (L\ γ\ u\ :\ ℝ≥0)\ (hg\ :\ LipschitzWith\ L\ g)}
\leanline{4}{(h\ :\ ℝ)\ (hh\ :\ 0\ ≤\ h)\ (initial\ :\ E\ ×\ E)\ (W\ :\ ℕ\ →\ ℝ\ →\ E)}
\leanline{4}{(hW\ :\ ∀\ n,\ Continuous\ (W\ n))\ :\ ℕ\ →\ E\ ×\ E}
\leanline{2}{|\ 0\ =>\ initial}
\leanline{2}{|\ n+1\ =>\ kineticSegmentFlow\ g\ L\ γ\ u\ hg\ h\ hh}
\leanline{6}{(kineticReferenceChain\ g\ L\ γ\ u\ hg\ h\ hh\ initial\ W\ hW\ n)\ (W\ n)\ (hW\ n)\ h}
\end{leancode}
\leansource{\texttt{kineticReferenceSegment} \quad KineticLangevin/\allowbreak{}CoupledSegments.lean, lines 113–116}
\begin{leancode}
\leanline{0}{def\ kineticReferenceSegment\ (g\ :\ E\ →\ E)\ (L\ γ\ u\ :\ ℝ≥0)\ (hg\ :\ LipschitzWith\ L\ g)}
\leanline{4}{(h\ :\ ℝ)\ (hh\ :\ 0\ ≤\ h)\ (initial\ :\ E\ ×\ E)\ (W\ :\ ℕ\ →\ ℝ\ →\ E)}
\leanline{4}{(hW\ :\ ∀\ n,\ Continuous\ (W\ n))\ (n\ :\ ℕ)\ :\ ℝ\ →\ E\ ×\ E\ :=}
\leanline{2}{kineticSegmentFlow\ g\ L\ γ\ u\ hg\ h\ hh\ (kineticReferenceChain\ g\ L\ γ\ u\ hg\ h\ hh\ initial\ W\ hW\ n)\ (W\ n)\ (hW\ n)}
\end{leancode}
\leansource{\texttt{exponentialInputChain} \quad KineticLangevin/\allowbreak{}CoupledSegments.lean, lines 119–123}
\begin{leancode}
\leanline{0}{def\ exponentialInputChain\ (g\ :\ E\ →\ E)\ (γ\ h\ :\ ℝ)\ (initial\ :\ E\ ×\ E)}
\leanline{4}{(B\ :\ ℕ\ →\ ℝ\ →\ E)\ :\ ℕ\ →\ E\ ×\ E}
\leanline{2}{|\ 0\ =>\ initial}
\leanline{2}{|\ n+1\ =>\ exponentialStep\ g\ γ\ h\ (exponentialInputChain\ g\ γ\ h\ initial\ B\ n)}
\leanline{6}{(B\ n\ h,∫\ s\ in\ 0..h,\ B\ n\ s)}
\end{leancode}
\leansource{\texttt{exponentialInputChain\_\allowbreak{}succ} \quad KineticLangevin/\allowbreak{}CoupledSegments.lean, lines 208–212}
\begin{leancode}
\leanline{0}{theorem\ exponentialInputChain\_succ\ (g\ :\ E\ →\ E)\ (γ\ h\ :\ ℝ)\ (initial\ :\ E\ ×\ E)}
\leanline{4}{(B\ :\ ℕ\ →\ ℝ\ →\ E)\ (n\ :\ ℕ)\ :}
\leanline{4}{exponentialInputChain\ g\ γ\ h\ initial\ B\ (n+1)\ =}
\leanline{6}{exponentialStep\ g\ γ\ h\ (exponentialInputChain\ g\ γ\ h\ initial\ B\ n)}
\leanline{8}{(B\ n\ h,∫\ s\ in\ 0..h,\ B\ n\ s)}
\end{leancode}
\leanlabel{How the Lean proof follows the argument}
\texttt{kinetic\allowbreak{}Reference\allowbreak{}Chain} names the mathematical object defined above. Its defining expression appears in the adjacent statement; subsequent identities may unfold this definition directly.




\texttt{kinetic\allowbreak{}Reference\allowbreak{}Segment} names the mathematical object defined above. Its defining expression appears in the adjacent statement; subsequent identities may unfold this definition directly.




\texttt{exponential\allowbreak{}Input\allowbreak{}Chain} names the mathematical object defined above. Its defining expression appears in the adjacent statement; subsequent identities may unfold this definition directly.




\texttt{exponential\allowbreak{}Input\allowbreak{}Chain\_\allowbreak{}succ} is a definitional equality: \texttt{rfl} closes it because both sides reduce to the same expression.
\end{leanbox}

\section{\texorpdfstring{The auxiliary exact path from the current numerical state}{The auxiliary exact path from the current numerical state}}


The auxiliary segment is a selected exact kinetic solution whose initial condition is the current exponential-integrator state. It receives exactly the same raw Brownian input as the reference segment on that iteration. The two specification theorems prove continuity, exact starting values, reference endpoint matching, and their shared-input integral equations. All these properties follow from the kinetic integral equation; no contraction or local-error inequality is inserted into the definitions.



At step \(n\), start an auxiliary exact segment at the numerical phase and drive it by the same input as the reference segment. It agrees initially with the numerical state, enabling the freezing-error estimate. Its discrepancy from the reference is controlled by exact-flow contraction.

\begin{leanbox}{Lean: \texttt{kineticAuxiliarySegment}, \texttt{kineticReferenceSegment\_\allowbreak{}spec}, \texttt{kineticAuxiliarySegment\_\allowbreak{}spec}}
\leansource{\texttt{kineticAuxiliarySegment} \quad KineticLangevin/\allowbreak{}CoupledSegments.lean, lines 126–129}
\begin{leancode}
\leanline{0}{def\ kineticAuxiliarySegment\ (g\ :\ E\ →\ E)\ (L\ γ\ :\ ℝ≥0)\ (hg\ :\ LipschitzWith\ L\ g)}
\leanline{4}{(h\ :\ ℝ)\ (hh\ :\ 0\ ≤\ h)\ (initial\ :\ E\ ×\ E)\ (W\ B\ :\ ℕ\ →\ ℝ\ →\ E)}
\leanline{4}{(hW\ :\ ∀\ n,\ Continuous\ (W\ n))\ (n\ :\ ℕ)\ :\ ℝ\ →\ E\ ×\ E\ :=}
\leanline{2}{kineticSegmentFlow\ g\ L\ γ\ 1\ hg\ h\ hh\ (exponentialInputChain\ g\ γ\ h\ initial\ B\ n)\ (W\ n)\ (hW\ n)}
\end{leancode}
\leansource{\texttt{kineticReferenceSegment\_\allowbreak{}spec} \quad KineticLangevin/\allowbreak{}CoupledSegments.lean, lines 132–149}
\begin{leancode}
\leanline{0}{theorem\ kineticReferenceSegment\_spec\ (g\ :\ E\ →\ E)\ (L\ γ\ u\ :\ ℝ≥0)\ (hg\ :\ LipschitzWith\ L\ g)}
\leanline{4}{(h\ :\ ℝ)\ (hh\ :\ 0\ ≤\ h)\ (initial\ :\ E\ ×\ E)\ (W\ :\ ℕ\ →\ ℝ\ →\ E)}
\leanline{4}{(hW\ :\ ∀\ n,\ Continuous\ (W\ n))\ (hW0\ :\ ∀\ n,\ W\ n\ 0\ =\ 0)\ (n\ :\ ℕ)\ :}
\leanline{4}{Continuous\ (kineticReferenceSegment\ g\ L\ γ\ u\ hg\ h\ hh\ initial\ W\ hW\ n)\ ∧}
\leanline{6}{kineticReferenceSegment\ g\ L\ γ\ u\ hg\ h\ hh\ initial\ W\ hW\ n\ 0\ =}
\leanline{8}{kineticReferenceChain\ g\ L\ γ\ u\ hg\ h\ hh\ initial\ W\ hW\ n\ ∧}
\leanline{6}{kineticReferenceSegment\ g\ L\ γ\ u\ hg\ h\ hh\ initial\ W\ hW\ n\ h\ =}
\leanline{8}{kineticReferenceChain\ g\ L\ γ\ u\ hg\ h\ hh\ initial\ W\ hW\ (n+1)\ ∧}
\leanline{6}{∀\ t\ ∈\ Icc\ 0\ h,}
\leanline{8}{kineticReferenceSegment\ g\ L\ γ\ u\ hg\ h\ hh\ initial\ W\ hW\ n\ t\ -}
\leanline{10}{kineticReferenceSegment\ g\ L\ γ\ u\ hg\ h\ hh\ initial\ W\ hW\ n\ 0\ =}
\leanline{10}{-(∫\ s\ in\ 0..t,\ kineticNegativeDrift\ g\ γ\ u}
\leanline{12}{(kineticReferenceSegment\ g\ L\ γ\ u\ hg\ h\ hh\ initial\ W\ hW\ n\ s))\ +\ (W\ n\ t,0)}
\end{leancode}
\leansource{\texttt{kineticAuxiliarySegment\_\allowbreak{}spec} \quad KineticLangevin/\allowbreak{}CoupledSegments.lean, lines 153–168}
\begin{leancode}
\leanline{0}{theorem\ kineticAuxiliarySegment\_spec\ (g\ :\ E\ →\ E)\ (L\ γ\ :\ ℝ≥0)\ (hg\ :\ LipschitzWith\ L\ g)}
\leanline{4}{(h\ :\ ℝ)\ (hh\ :\ 0\ ≤\ h)\ (initial\ :\ E\ ×\ E)\ (W\ B\ :\ ℕ\ →\ ℝ\ →\ E)}
\leanline{4}{(hW\ :\ ∀\ n,\ Continuous\ (W\ n))\ (hW0\ :\ ∀\ n,\ W\ n\ 0\ =\ 0)\ (n\ :\ ℕ)\ :}
\leanline{4}{Continuous\ (kineticAuxiliarySegment\ g\ L\ γ\ hg\ h\ hh\ initial\ W\ B\ hW\ n)\ ∧}
\leanline{6}{kineticAuxiliarySegment\ g\ L\ γ\ hg\ h\ hh\ initial\ W\ B\ hW\ n\ 0\ =}
\leanline{8}{exponentialInputChain\ g\ γ\ h\ initial\ B\ n\ ∧}
\leanline{6}{∀\ t\ ∈\ Icc\ 0\ h,}
\leanline{8}{kineticAuxiliarySegment\ g\ L\ γ\ hg\ h\ hh\ initial\ W\ B\ hW\ n\ t\ -}
\leanline{10}{kineticAuxiliarySegment\ g\ L\ γ\ hg\ h\ hh\ initial\ W\ B\ hW\ n\ 0\ =}
\leanline{10}{-(∫\ s\ in\ 0..t,\ kineticNegativeDrift\ g\ γ\ 1}
\leanline{12}{(kineticAuxiliarySegment\ g\ L\ γ\ hg\ h\ hh\ initial\ W\ B\ hW\ n\ s))\ +\ (W\ n\ t,0)}
\end{leancode}
\leanlabel{How the Lean proof follows the argument}
\texttt{kinetic\allowbreak{}Auxiliary\allowbreak{}Segment} names the mathematical object defined above. Its defining expression appears in the adjacent statement; subsequent identities may unfold this definition directly.




\texttt{kinetic\allowbreak{}Reference\allowbreak{}Segment\_\allowbreak{}spec} gives the formal counterpart of this step. The main intermediate facts are \texttt{hs} in that order. Named proof applications include \texttt{hs.\allowbreak{}2.\allowbreak{}2}; the adjacent proof shows their hypotheses and order of use.




\texttt{kinetic\allowbreak{}Auxiliary\allowbreak{}Segment\_\allowbreak{}spec} gives the formal counterpart of this step. The main intermediate facts are \texttt{hs} in that order. Named proof applications include \texttt{hs.\allowbreak{}2.\allowbreak{}2}; the adjacent proof shows their hypotheses and order of use.
\end{leanbox}

\section{\texorpdfstring{Discharge the local-error proof's kinematic and mild requirements}{Discharge the local-error proof's kinematic and mild requirements}}


The constructed auxiliary path has its positional primitive identity by the general kinetic projection theorem. Its exponentially weighted force formula follows by restricting its integral equation to any smaller time \(t\le h\), restricting the filtered-input balance to the same interval, and applying the generic mild theorem. The argument explicitly checks the interval inclusions and simplifies the inverse mass \(u=1\). The remaining stochastic work concerns the law and stationarity of the reference; the coupling segments and their integral identities have been constructed.



The local-error argument needs continuity, the mild equation, and the integrated-velocity position equation. Each follows from the selected solution’s specification. These hypotheses are therefore verified when the constructed segments are inserted into the one-step coupling theorem.

\begin{leanbox}{Lean: \texttt{kineticAuxiliarySegment\_\allowbreak{}kinematic}, \texttt{kineticAuxiliarySegment\_\allowbreak{}mild}}
\leansource{\texttt{kineticAuxiliarySegment\_\allowbreak{}kinematic} \quad KineticLangevin/\allowbreak{}CoupledSegments.lean, lines 171–179}
\begin{leancode}
\leanline{0}{theorem\ kineticAuxiliarySegment\_kinematic\ (g\ :\ E\ →\ E)\ (L\ γ\ :\ ℝ≥0)\ (hg\ :\ LipschitzWith\ L\ g)}
\leanline{4}{(h\ :\ ℝ)\ (hh\ :\ 0\ ≤\ h)\ (initial\ :\ E\ ×\ E)\ (W\ B\ :\ ℕ\ →\ ℝ\ →\ E)}
\leanline{4}{(hW\ :\ ∀\ n,\ Continuous\ (W\ n))\ (hW0\ :\ ∀\ n,\ W\ n\ 0\ =\ 0)\ (n\ :\ ℕ)}
\leanline{4}{(t\ :\ ℝ)\ (ht\ :\ t\ ∈\ Icc\ 0\ h)\ :}
\leanline{4}{(kineticAuxiliarySegment\ g\ L\ γ\ hg\ h\ hh\ initial\ W\ B\ hW\ n\ t).2\ -}
\leanline{6}{(kineticAuxiliarySegment\ g\ L\ γ\ hg\ h\ hh\ initial\ W\ B\ hW\ n\ 0).2\ =}
\leanline{6}{∫\ s\ in\ 0..t,\ (kineticAuxiliarySegment\ g\ L\ γ\ hg\ h\ hh\ initial\ W\ B\ hW\ n\ s).1}
\end{leancode}
\leansource{\texttt{kineticAuxiliarySegment\_\allowbreak{}mild} \quad KineticLangevin/\allowbreak{}CoupledSegments.lean, lines 182–205}
\begin{leancode}
\leanline{0}{theorem\ kineticAuxiliarySegment\_mild\ (g\ :\ E\ →\ E)\ (L\ γ\ :\ ℝ≥0)\ (hg\ :\ LipschitzWith\ L\ g)}
\leanline{4}{(h\ :\ ℝ)\ (hh\ :\ 0\ ≤\ h)\ (initial\ :\ E\ ×\ E)\ (W\ B\ :\ ℕ\ →\ ℝ\ →\ E)}
\leanline{4}{(hW\ :\ ∀\ n,\ Continuous\ (W\ n))\ (hW0\ :\ ∀\ n,\ W\ n\ 0\ =\ 0)}
\leanline{4}{(hB\ :\ ∀\ n,\ Continuous\ (B\ n))}
\leanline{4}{(hnoise\ :\ ∀\ n\ t,\ t\ ∈\ Icc\ 0\ h\ →\ B\ n\ t\ =\ W\ n\ t\ -\ (γ\ :\ ℝ)\ •\ ∫\ s\ in\ 0..t,\ B\ n\ s)}
\leanline{4}{(n\ :\ ℕ)\ (t\ :\ ℝ)\ (ht\ :\ t\ ∈\ Icc\ 0\ h)\ :}
\leanline{4}{(kineticAuxiliarySegment\ g\ L\ γ\ hg\ h\ hh\ initial\ W\ B\ hW\ n\ t).1\ =}
\leanline{6}{psi\ γ\ 0\ t\ •\ (kineticAuxiliarySegment\ g\ L\ γ\ hg\ h\ hh\ initial\ W\ B\ hW\ n\ 0).1\ -}
\leanline{6}{(∫\ s\ in\ 0..t,\ Real.exp\ (-(γ\ :\ ℝ)*(t-s))\ •}
\leanline{8}{g\ (kineticAuxiliarySegment\ g\ L\ γ\ hg\ h\ hh\ initial\ W\ B\ hW\ n\ s).2)\ +\ B\ n\ t}
\end{leancode}
\leanlabel{How the Lean proof follows the argument}
\texttt{kinetic\allowbreak{}Auxiliary\allowbreak{}Segment\_\allowbreak{}kinematic} gives the formal counterpart of this step. The main intermediate facts are \texttt{hs} in that order. Named proof applications include \texttt{kinetic\_\allowbreak{}integral\_\allowbreak{}kinematic}; the adjacent proof shows their hypotheses and order of use.




\texttt{kinetic\allowbreak{}Auxiliary\allowbreak{}Segment\_\allowbreak{}mild} gives the formal counterpart of this step. The main intermediate facts are \texttt{hs}, \texttt{heq}, \texttt{hb} in that order. Named proof applications include \texttt{hs.\allowbreak{}2.\allowbreak{}2}, \texttt{hnoise}, \texttt{kinetic\_\allowbreak{}integral\_\allowbreak{}mild}; the adjacent proof shows their hypotheses and order of use.
\end{leanbox}

\end{document}
